\documentclass[10pt,a4paper]{article}
\usepackage{amsmath,amssymb,amsthm,graphicx}
\newcommand{\bits}{\{0,1\}}
\newcommand{\DT}{\operatorname{DT}}
\newcommand{\DTSize}{\operatorname{DTSize}}
\newcommand{\CC}{\operatorname C}
\newcommand{\comp}{\diamond}
\newcommand{\OR}{\operatorname{OR}}
\newcommand{\AND}{\operatorname{AND}}
\newcommand{\XOR}{\operatorname{XOR}}
\newenvironment{originalresult}[1]{\par\medskip\noindent\textbf{#1.}\itshape}{\par\medskip}
\title{Complete gadget dichotomy for randomized communication lifting}
\author{Yaroslav Alekseev, Yuval Filmus, Alexander V. Smal}
\date{doi:10.1007/s00037-025-00276-5; CC BY 4.0}
\begin{document}
\maketitle
\noindent\textit{Editorial note.} Actual complete CC BY4.0 journal version; faithful editable transcription. Full author Theorem6.1 reduction, Lemma6.3 matrix proof and Lemma6.8 protocol are retained as uncounted support for Theorem6.7. Standard disjointness hardness and equality protocol are precise cited prerequisites. Mathematical body uses the actual non-reducibility symbols from source, without adopting text-extraction glyph loss. No author mathematics is repaired.\par\medskip
\noindent\textit{Allowed standard prerequisites (editorial).}\begin{itemize}
\item Randomized unique set-disjointness lower bound. R(PromiseDisj\_m)=Omega(m). Razborov1992; exact use retained in full author Theorem6.1 proof
\item Private-coin randomized equality protocol. Kushilevitz and Nisan1996; exact O(log n) use on EQ\_nm retained
\end{itemize}\medskip
\section*{Block composition (author source p.2)}
For $f:\bits^n\to V$ and $g:\bits^m\to\bits$, a (block-)composition $f\comp g:\bits^{n\times m}\to V$ is defined by
\[(f\comp g)(z_1,z_2,\ldots,z_n):=f(g(z_1),g(z_2),\ldots,g(z_n)),\]
where each $z_i\in\bits^m$.
\section*{Prerequisites (author source pp.10--12)}
All the logarithms are base 2. For $g:\bits^k\to\bits$, we define a function $g^n:\bits^{n\times k}\to\bits^n$ such that
\[g^n(x_1,x_2,\ldots,x_n):=(g(x_1),g(x_2),\ldots,g(x_n)).\]
For a function $f:\bits^n\to\bits$, an input $x\in\bits^n$, and a partial assignment $\sigma\in\{0,1,*\}^n$ that agrees with $x$ on all non-$*$ coordinates, we use $f|_\sigma$ to denote a restriction of $f$ to $\sigma$ and $x|_\sigma$ to denote a projection of $x$ to the $*$-coordinates of $\sigma$.
\paragraph{Decision tree complexity.} A decision tree is a rooted binary tree with internal nodes labeled by input variables (representing queries), edges labeled with 0 and 1, and leaves labeled by some values. Decision tree evaluation is defined in a natural way: given an assignment for the input variables, we traverse the tree starting from the root and then sequentially choose every next edge according to the given assignments until we reach some leaf. The label of this leaf is the result of the evaluation. A decision tree computes some function $f$ if for all possible assignments to the input variables, the decision tree evaluates to the value of the function. For a function $f$, $\DT(f)$ denotes the minimal depth of a decision tree computing $f$, and $\DTSize(f)$ denotes the minimal number of leaves in a decision tree computing $f$.
\paragraph{Certificate complexity.} A certificate complexity is a non-deterministic analogue of the decision tree complexity. A certificate for an input $x\in\bits^n$ to a function $f$ is a set $S\subseteq[n]$ of indices such that $f$ restricted to all inputs that match $x$ on $S$ is constant, i.e., $f(y)=f(x)$ whenever $y|_S=x|_S$. The certificate complexity $\CC(f,x)$ of $f$ at input $x$ is the size of the smallest certificate for $x$. Finally, the certificate complexity of function $f$ is defined as $\CC(f):=\max_{x\in\bits^n}\CC(f,x)$.

\section*{Block sensitivity and randomized communication (author source pp.12--13)}
A block $B\subseteq[n]$ is sensitive for a function $f$ at $x$ iff $f(x)\ne f(x^B)$. The block sensitivity $\operatorname{bs}(f,x)$ of $f$ at $x$ is the maximum number of disjoint sensitive blocks for $f$ at $x$. The block sensitivity $\operatorname{bs}(f)$ of $f$ is the maximum sensitivity $\operatorname{bs}(f,x)$ over all points $x\in\bits^n$.
We use the notation $x^B$ for $x\in\bits^n$ and $B\subseteq[n]$ to denote $x$ with all the coordinates in $B$ flipped.
In a randomized communication game, Alice and Bob have access to an unlimited amount of random bits, and they can use it to decide which bit to send next. We say that a (private coin) randomized communication protocol $\Pi$ computes a function $f:X\times Y\to Z$ with error $\varepsilon$ if
\[\Pr_{r_A,r_B}[\Pi(x,y,r_A,r_B)=f(x,y)]\ge1-\varepsilon,\qquad\forall x\in X,y\in Y,\]
where $r_A$ and $r_B$ are the strings of random bits used by Alice and Bob, respectively. The cost of a randomized protocol is the maximum number of bits that can be sent. We denote by $R_\varepsilon(f)$ the minimal cost of a private coin randomized communication protocol that computes $f$ with error $\varepsilon$. The randomized communication complexity of $f$ is defined as $R(f):=R_{1/3}(f)$.
\section*{6.1 Reductions in communication complexity (author source p.35)}
Consider a pair of two-party functions $f_i:X_i\times Y_i\to\bits$ for $i=1,2$. We say that $f_1$ reduces to $f_2$, denoted $f_1\le f_2$, if the communication matrix of $f_1$ is a submatrix of the communication matrix of $f_2$. Equivalently, $f_1\le f_2$ iff there exist one-to-one mappings $\pi_A$ and $\pi_B$ such that
\[f_1(x,y)=f_2(\pi_A(x),\pi_B(y)),\qquad\forall(x,y)\in X_1\times Y_1.\]

\begin{originalresult}{Theorem 6.1 (Zhang; author source p.36)}
If a two-party gadget $g:X\times Y\to\bits$ satisfies $\AND,\OR\le g$, then for every function $f:\bits^n\to Q$, the function $f\comp g$ has (constant error) randomized communication complexity $\Omega(\operatorname{bs}(f))$.
\end{originalresult}
For completeness, we present the proof that is a reduction from unique set disjointness.
\begin{originalresult}{Definition 6.2}
The disjointness function $\operatorname{Disj}_n:\bits^n\times\bits^n$ is defined as follows:
\[\operatorname{Disj}_n(x,y):=\bigvee_{i=1}^n x_i\wedge y_i.\]
Unique set disjointness, $\operatorname{PromiseDisj}_n$, is a promise version of disjointness with the promise that there is always at most one $i$ such that $x_i\wedge y_i=1$.
\end{originalresult}
\begin{proof}[Proof of Theorem 6.1]
Razborov (1992) showed that $R(\operatorname{PromiseDisj}_m)=\Omega(m)$. We will show a reduction from this problem to computing $f\comp g$. To illustrate the proof, let us start with the case in which $f(0)=0$ and $f(e_i)=1$ for all $i\in[n]$, where $e_i$ is the $i$'th unit vector. Then
\[\operatorname{PromiseDisj}_n(x,y)=f(x_1\wedge y_1,\ldots,x_n\wedge y_n).\]
Since $\AND\le g$, we can solve unique set disjointness using a protocol for $f\comp g$. Therefore, $f\comp g$ has communication complexity $\Omega(n)$. Note that this case only required $\AND\le g$; we need $\OR\le g$ in order to handle inputs changing from 1 to 0.
To handle the general case, let inputs $X^\wedge(0),X^\wedge(1)\in X$ and $Y^\wedge(0),Y^\wedge(1)\in Y$ witness $\AND\le g$, and let $X^\vee(0),X^\vee(1)\in X$ and $Y^\vee(0),Y^\vee(1)\in Y$ witness $\OR\le g$.
Let $z$ be a point witnessing $\operatorname{bs}(f)$, and let $B_1,\ldots,B_m$ be the corresponding blocks, where $m=\operatorname{bs}(f)$. Given inputs $x,y$ to unique set disjointness, we construct the following inputs $X,Y$ to $f\comp g$:
\begin{itemize}
\item If $i$ does not belong to any of the blocks, then we fix $X_i,Y_i$ so that $g(X_i,Y_i)=z_i$.
\item If $i\in B_j$ and $z_i=0$, put $X_i=X^\wedge(x_j),Y_i=Y^\wedge(y_j)$.
\item If $i\in B_j$ and $z_i=1$, put $X_i=X^\vee(\bar x_j),Y_i=Y^\vee(\bar y_j)$.
\end{itemize}
The reader can check that if $x,y$ are disjoint, then applying $g$ to $X,Y$ yields $z$, and if they intersect only at $j$, then it yields $z\oplus B_j$. Hence, $f\comp g$ can be used to compute unique set disjointness of length $m$, and consequently, we obtain a lower bound of $\Omega(m)$.
\end{proof}

\section*{Block-diagonal gadgets (author source pp.37--38)}
For the classification of gadgets, we will need the following lemma that shows that if $\OR\nleq g$ then the communication matrix of $g$ is block diagonal.
\begin{originalresult}{Lemma 6.3}
If $g:X\times Y\to\bits$ satisfies that $\OR\nleq g$ then there are partitions of $X$ and $Y$ into disjoint sets
\[X=X_0\sqcup X_1\sqcup\cdots\sqcup X_k\quad\text{and }Y=Y_0\sqcup Y_1\sqcup\cdots\sqcup Y_k,\]
such that
\begin{itemize}
\item If $x\in X_i,y\in Y_j$, where $i\ne j$, or $i=0$, or $j=0$, then $g(x,y)=0$.
\item If $x\in X_i,y\in Y_i$, where $i\ne0$, then $g(x,y)=1$.
\end{itemize}
\end{originalresult}
\begin{proof}
We can construct such a partition step by step. Consider any point $(x_0,y_0)$ such that $g(x_0,y_0)=1$. If we don't have such a point, then we can take $X_0=X$ and $Y_0=Y$. Otherwise, we consider all points $y\in Y$ such that $g(x_0,y)=1$. Denote the set of all such points as $Y'$. Now, we can consider all $x\in X$ such that $g(x,y_0)=1$. Denote the set of all such points as $X'$. Observe that every point $(x,y)\in X'\times Y'$ satisfies $g(x,y)=1$, otherwise we can project our gadget to points $(x_0,y_0),(x_0,y),(x,y_0),(x,y)$ and get OR.
Now, observe that for any $(x'',y')\in(X\setminus X')\times Y'$ we have $g(x'',y_0)=0$ by construction, and thus $g(x'',y')=0$, since otherwise we can project the gadget to points $(x_0,y'),(x_0,y_0),(x'',y'),(x'',y_0)$ and get OR. By analogy, for any $(x',y'')\in X'\times(Y\setminus Y')$ we have $g(x',y'')=0$.
This means that $X'\times Y'$ forms a valid block of our partition, so we can use the same procedure to construct a block of $(X\setminus X')\times(Y\setminus Y')$. In the end, if we get sets $X''$ and $Y''$ such that any $(x,y)\in X''\times Y''$ satisfies $g(x,y)=0$, we take $X_0=X''$ and $Y_0=Y''$.
\end{proof}

\section*{6.2 Gadget classification for randomized communication complexity (author source pp.39--40)}
In this section, we prove the following dichotomy result for the randomized case.

\begin{originalresult}{Theorem 6.7}
For every gadget $g$, one of the following cases holds:
\begin{enumerate}
\item[(i)] There is an infinite family of functions $f_n$ with $\operatorname{bs}(f_n)=\Omega(n)$ and $R(f_n\comp g)=O(\log n)$.
\item[(ii)] For every function $f$, we have $R(f\comp g)=\Omega(\operatorname{bs}(f)^{\Omega(1)})$.
\end{enumerate}
\end{originalresult}

First, we need to prove the following lemma.
\begin{originalresult}{Lemma 6.8}
Suppose that gadgets $g,h:X\times Y\to\bits$ are such that $\OR\nleq g$ and $\AND\nleq h$. Then
\[R(\AND_n\comp g)=O(\log n),\qquad R(\OR_n\comp h)=O(\log n).\]
\end{originalresult}
\begin{proof}
The bounds for $g$ and $h$ are symmetrical, so we prove only the first one. Consider the decomposition $X=X_0\sqcup X_1\sqcup\cdots\sqcup X_k$ and $Y=Y_0\sqcup Y_1\sqcup\cdots\sqcup Y_k$ provided by Lemma 6.3. If one of the blocks $g_i$ was given an input from $X_0$ or $Y_0$, then one of the players can share this information with the other player and finish the communication. Otherwise, they indicate that they do not have such blocks. After this step, they are restricted to inputs from $X\setminus X_0$ and $Y\setminus Y_0$. Since any two inputs from the same block $X_i$ or $Y_i$ are indistinguishable for our gadget, they need to solve $\AND_n\comp\operatorname{EQ}_m$, where $\operatorname{EQ}_m$ is the equality function for $m$-bit inputs (i.e., $\operatorname{EQ}_m(x,y)=1\iff x=y$) for $m=\lceil\log k\rceil$. This is exactly an instance of $\operatorname{EQ}_{nm}$, which can be solved with $O(\log n)$ bits of communication (note that $m$ is a constant) by a randomized communication protocol (see Kushilevitz \& Nisan (1996)).
\end{proof}
Now we are ready to prove the dichotomy result.
\begin{proof}[Theorem 6.7]
If $\AND,\OR\le g$ then Theorem 6.1 shows that $R(f\comp g)=\Omega(\operatorname{bs}(f)^{\Omega(1)})$, so we get the second case.
Otherwise, $\OR\nleq g$ or $\AND\nleq g$. We apply Lemma 6.8: if $\OR\nleq g$ we take $f_n=\AND_n$, otherwise we take $f_n=\OR_n$. This gives us the first case of the theorem.
\end{proof}

\begin{thebibliography}{9}
\bibitem{zhang} Shengyu Zhang (2009). On the Tightness of the Buhrman--Cleve--Wigderson Simulation. In Proceedings of the 20th International Symposium on Algorithms and Computation,ISAAC'09,434--440. Springer. doi:10.1007/978-3-642-10631-6\_45.
\bibitem{razborov} Alexander Razborov (1992). On the distributional complexity of disjointness. Theoretical Computer Science106(2),385--390.
\bibitem{kn} Eyal Kushilevitz and Noam Nisan (1996). Communication Complexity. Cambridge University Press.
\bibitem{hatami} Hamed Hatami, Kaave Hosseini and Shachar Lovett (2018). Structure of protocols for XOR functions. SIAM Journal on Computing47(1),208--217. doi:10.1137/17M1136869.
\bibitem{zhangshi} Zhiqiang Zhang and Yaoyun Shi (2010). On the parity complexity measures of Boolean functions. Theoretical Computer Science411(26),2612--2618.
\end{thebibliography}
\end{document}
